%% ma: L10-B12-0003
%% lop: 10
%% chuong: 5
%% bai: 12
%% dang: 10.12.2
%% loai: TL
%% muc: TH
%% dap_an: "-"
%% nguon: "chuyen-de-cac-so-dac-trung-cua-mau-so-lieu-khong-ghep-nhom-toan-10-chuong-trinh-moi.pdf (GV. Phan Nhật Linh) (P01), Bài 01 Bài tập tự luận 3"
%% kien_thuc: "Lớp 10 Bài 12 (số gần đúng, sai số tuyệt đối, độ chính xác, sai số tương đối, quy tròn số)"
%% trang_thai: chinh-thuc
%% ghi_chu: "Tự duyệt theo yêu cầu giáo viên 22/10/2026."
\begin{ex}
Các nhà toán học cổ đại Trung Quốc đã dùng phân số $\dfrac{22}{7}$ để xấp xỉ số $\pi$. Hãy đánh giá sai số tuyệt đối của giá trị gần đúng này, biết $3{,}1415<\pi<3{,}1416$.
\loigiai{$\dfrac{22}7=3{,}142857\ldots<3{,}1429$ và $\pi>3{,}1415$ nên $\pi<\dfrac{22}7$.

Sai số tuyệt đối $\Delta=\left|\pi-\dfrac{22}7\right|=\dfrac{22}7-\pi<3{,}1429-3{,}1415=0{,}0014$.}
\end{ex}
